\begin{definition}[Normal PEPS blocking hypotheses]%
    \label{def:peps_normalPEPSBlockingHypotheses}
    \lean{TNLean.PEPS.NormalPEPSBlockingHypotheses}
    \lean{TNLean.PEPS.NormalPEPSBlockingHypotheses.injective_chain_at_edge}
    \lean{TNLean.PEPS.%
        NormalPEPSBlockingHypotheses.endpoint_disjoint_cover_at_edge}
    \lean{TNLean.PEPS.%
        NormalPEPSBlockingHypotheses.endpoint_injective_disjoint_cover_at_edge}
    \lean{TNLean.PEPS.NormalPEPSBlockingHypotheses.mem_withSite_iff}
    \leanok
    \uses{def:peps_normalEdgeBlockingHypotheses,
        def:peps_normalOneSiteSeparationHypotheses}
    The general normal theorem assumes both the edge-centred blockings into
    three injective regions and the one-site comparison regions with injective
    complements. These two assumptions are collected as the normal PEPS
    blocking hypotheses; the same hypotheses give the corresponding edge and
    one-site consequences without restating the separate edge and one-site
    assumptions.
    \begin{center}
        \TNPEPSNormalBlockingHypotheses
    \end{center}
\end{definition}

\begin{theorem}[Square-lattice normal blocking hypotheses from margin covers]%
    \label{thm:peps_normalSquarePEPSBlockingHypotheses_of_marginCovers}
    \lean{TNLean.PEPS.normalSquarePEPSBlockingHypotheses_of_marginCovers}
    \lean{TNLean.PEPS.normalSquarePEPSBlockingHypotheses_edgeBlocking_of_marginCovers}
    \lean{TNLean.PEPS.normalSquarePEPSBlockingHypotheses_blockingData_of_marginCovers}
    \lean{TNLean.PEPS.%
        normalSquarePEPSBlockingHypotheses_oneSiteSeparation_of_marginCovers}
    \lean{TNLean.PEPS.%
        normalSquarePEPSBlockingHypotheses_injective_chain_of_marginCovers}
    \lean{TNLean.PEPS.%
        normalSquarePEPSBlockingHypotheses_endpoint_disjoint_cover_of_marginCovers}
    \lean{TNLean.PEPS.%
        normalSquarePEPSBlockingHypotheses_endpoint_injective_disjoint_cover_of_marginCovers}
    \lean{TNLean.PEPS.%
        normalSquarePEPSBlockingHypotheses_mem_withSite_iff_of_marginCovers}
    \leanok
    \uses{thm:peps_normalSquareTranslatedEdgeBlockingHypotheses_of_windows,
        def:peps_normalPEPSBlockingHypotheses,
        def:peps_normalOneSiteSeparationHypotheses}
    Given rectangular injectivity hypotheses, finite union closure, oriented
    margin-and-cover hypotheses at every edge, and one-site separation, the finite
    square lattice satisfies the general normal blocking hypotheses. The
    edge-blocking hypotheses are exactly the margin-cover construction, the
    one-edge datum at an edge is the datum supplied there, and the one-site
    separation is exactly the supplied one-site separation. Hence these hypotheses
    give the three injective regions, the endpoint-disjoint-cover facts at
    each edge, and the one-site membership formula as consequences of the same
    normal blocking hypotheses.
\end{theorem}
\begin{proof}\leanok
    Combine the edge-blocking hypotheses constructed from the margin-cover
    hypotheses with the given one-site separation hypotheses.
\end{proof}

\begin{theorem}[Fundamental Theorem for normal PEPS]%
    \label{thm:fundamentalTheorem_normalPEPS}
    \notready
    \uses{thm:peps_normalEdgeBlockingHypotheses_injectiveChain,
        def:peps_normalPEPSBlockingHypotheses,
        thm:peps_oneVertexComplementComparison,
        def:GaugeEquivPEPS}
    Suppose two normal PEPS generate the same state, can be blocked into
    three-site injective chains around every edge, and for every site admit
    two injective regions with injective complements that differ exactly at
    that site. Then their defining tensors are related by a local gauge, and
    the gauges are unique up to the scalar freedom allowed by the graph. This
    is the general normal-PEPS theorem stated after
    \cite[Theorem~3]{Molnar2018NormalPEPS}.
    A connected-graph form, with the blocking hypothesis stated as the
    source's proof uses it (one blocking shared by the two states, each
    distinguished edge the entire bond of its chain), is
    Theorem~\ref{thm:fundamentalTheorem_normalPEPS_connected}; its per-edge
    uniqueness clause is
    Theorem~\ref{thm:fundamentalTheorem_normalPEPS_gauge_unique}.
    The source's remark following the theorem, that for a translationally
    invariant system the gauges are also translationally invariant when the
    proportionality constants are not absorbed into them, is realized on the
    torus by
    Theorem~\ref{thm:fundamentalTheorem_normalTorusPEPS_unconditional}: its
    gauge family is carried onto every edge by the translations from one
    reference matrix per orientation class, and its scalar $\lambda$ is kept
    separate rather than absorbed.
\end{theorem}

\begin{definition}[Doubly injective regions]%
    \label{def:peps_regionInjectivityDataPair}
    \lean{TNLean.PEPS.regionInjectivityDataPair}
    \leanok
    \uses{def:peps_regionInjectivityData}
    For two region-injectivity assignments on the same vertex set, a region is
    doubly injective when it is injective for each of the two.  The general
    normal theorem blocks two states by one shared geometry, and its final
    comparison contracts both tensors over the same regions, so its blocking
    hypotheses are instantiated at the doubly injective assignment of the two
    tensors.
\end{definition}

\begin{theorem}[Equal bond dimensions from the normal blocking hypotheses]%
    \label{thm:peps_bondDim_eq_of_normalBlocking}
    \lean{TNLean.PEPS.bondDim_eq_of_normalBlocking}
    \leanok
    \uses{def:peps_normalPEPSBlockingHypotheses,
        def:peps_regionInjectivityDataPair}
    Let $A$ and $B$ be PEPS tensors on a finite simple graph with the same
    positive physical dimension and positive bond dimensions, generating the
    same state and satisfying the normal blocking hypotheses at the doubly
    injective assignment of the two tensors, with each distinguished edge the
    only edge of the graph joining its two adjacent blocks.  Then the bond
    dimensions of $A$ and $B$ agree on every edge.
    In the source the equality is part of the isomorphism lemma of
    \cite[Section~3]{Molnar2018NormalPEPS}: the insertion correspondence
    $X\mapsto Y$ on the chosen bond is an algebra isomorphism between the two
    full bond matrix algebras, so the bond dimensions on the two sides agree;
    the normal case inherits the argument through the blocked three-site
    chains.  The single-crossing condition is the source's blocking around an
    edge made explicit: the chosen edge is the entire bond between the first
    two parties of the blocked chain.
    \begin{proof}
        \leanok
        \uses{thm:peps_edgeBlockedThreeSiteInjective,
            thm:peps_edgeBlockedInsertionAlgebraIsomorphism}
        Fix an edge.  Blocking both networks into the shared three regions
        around it produces two three-site chains on the three-vertex complete
        graph whose super-sites carry the blocked-region weight families,
        injective by hypothesis.  Each coarse state coefficient is a positive
        merge-collapse constant, a product of bond dimensions over the edges
        not incident to each region, times the corresponding original state
        coefficient, so the two coarse states are proportional with ratio the
        two constants.  Multiplying one super-site of each chain by the other
        chain's constant yields chains with equal states, unchanged bond
        dimensions, and the super-site injectivities preserved.  The matched
        matrix insertions on the distinguished coarse bond of the two scaled
        chains form an algebra isomorphism between the two full bond matrix
        algebras
        (Theorem~\ref{thm:peps_edgeBlockedInsertionAlgebraIsomorphism}), and
        an isomorphism of full matrix algebras forces equal matrix sizes.
        Under the single-crossing hypothesis the coarse bond carries the
        original bond dimension of the distinguished edge for each tensor, so
        the two original bond dimensions agree.
    \end{proof}
\end{theorem}

\begin{theorem}[Fundamental Theorem for normal PEPS on a connected graph]%
    \label{thm:fundamentalTheorem_normalPEPS_connected}
    \lean{TNLean.PEPS.fundamentalTheorem_normalPEPS}
    \leanok
    % Reading note: this is the general normal theorem stated after Theorem 3
    % of \cite{Molnar2018NormalPEPS} on a connected graph.  The shared
    % blockings and the single-crossing hypothesis are the source's blocking
    % hypothesis made explicit, not restrictions; both readings are recorded
    % in \path{docs/paper-gaps/peps_normal_ft_section3_route.tex}.
    \uses{def:peps_normalPEPSBlockingHypotheses,
        def:peps_regionInjectivityDataPair, def:GaugeEquivPEPS}
    Let $A$ and $B$ be PEPS tensors on a connected finite simple graph with
    the same positive physical dimension and positive bond dimensions,
    generating the same state.  Suppose the normal blocking
    hypotheses hold at the doubly injective assignment of the two tensors:
    around every edge the network blocks into three doubly injective regions
    with the edge endpoints in the first two, and every site admits two doubly
    injective regions, with doubly injective complements, differing exactly at
    that site.  Suppose moreover that at every edge the distinguished edge is
    the only edge of the graph joining its two adjacent blocks.  Then the
    defining tensors are related by a local gauge: there are invertible
    matrices, one on each edge, whose endpoint actions transform $A_v$ into
    $B_v$ at every vertex, with no leftover scalar.
    This is the general normal-PEPS theorem stated after
    \cite[Theorem~3]{Molnar2018NormalPEPS}
    (Theorem~\ref{thm:fundamentalTheorem_normalPEPS}), with the source's
    blocking hypothesis stated explicitly as its proof uses it, in two ways,
    neither of which restricts the source statement.  First, the blockings
    are shared between the two states, each region injective for both
    tensors: the source blocks both states by one blocking (``\emph{they}
    can be blocked''), its isomorphism lemma takes the two blocked chains
    over the same tripartition with all six blocks injective, and its final
    comparison contracts both tensors over the same one-site-different
    regions, whose injectivity for both tensors its comparison lemma
    requires.  Second, the single-crossing hypothesis is the explicit
    statement of the source's blocking \emph{around} an edge: the chosen edge
    is the entire bond between the first two parties of the three-site chain.
    In the source, blocking around an edge groups all vertices except the two
    endpoints of the edge, so the edge is the bond between the first two
    parties; the regions in the proof of
    \cite[Theorem~3]{Molnar2018NormalPEPS} enlarge those parties to injective
    blocks meeting only along the distinguished edge, and the gauge of the
    isomorphism lemma lives on that edge of the original network. The
    equality of the bond dimensions is not assumed; it is derived from the
    blocking hypotheses
    (Theorem~\ref{thm:peps_bondDim_eq_of_normalBlocking}), as in the source's
    isomorphism lemma. The
    positive bond dimensions and the connectivity of the graph are the
    counterexample-backed corrections carried over from the injective
    Fundamental Theorem (Theorem~\ref{thm:fundamentalTheorem_PEPS}); the
    absorption of the per-vertex constants into the gauges is exactly the step
    that requires connectivity.
    \begin{proof}
        \leanok
        \uses{thm:peps_bondDim_eq_of_normalBlocking,
            lem:peps_injectiveRegion_union_pos,
            thm:peps_region_edge_gauge_of_coeff_transfer,
            thm:gaugeConsistency}
        The bond dimensions of the two tensors agree on every edge by
        Theorem~\ref{thm:peps_bondDim_eq_of_normalBlocking}.
        The union closure of injective regions for either tensor follows from
        the positive bond dimensions and
        Lemma~\ref{lem:peps_injectiveRegion_union_pos}.  At every edge the
        shared blocking datum restricts to a datum for each tensor over the
        same three regions, so the per-edge gauge of the region insertion
        transfer (Theorem~\ref{thm:peps_region_edge_gauge_of_coeff_transfer})
        and its orientation-adapted absorption give a family of invertible
        matrices realizing the absorbed inserted-coefficient equality at every
        edge.  A gauge preserves blocked-region linear independence, so at
        every site the one-site comparison applies to the gauge-absorbed
        tensor: comparing the smaller region with its one-site completion
        yields two proportionality scalars whose ratio is a nonzero scalar
        $\lambda_v$ with $A_v$ equal to $\lambda_v$ times the gauge action on
        $B$ at $v$.  When the smaller region is empty its block is the count
        of global virtual configurations, and when the completion is the full
        vertex set its block is the closed state coefficient; in these
        degenerate cases the proportionality holds with scalar one, and every
        other comparison region has a boundary edge by connectivity.  The
        per-vertex relations against the nonzero closed state force
        $\prod_v \lambda_v = 1$, so on the connected graph a spanning-tree
        construction realizes the reciprocals of the scalars as oriented
        incidence products of edge scalars, which fold into the gauges as in
        the absorption step of the injective Fundamental Theorem
        (Theorem~\ref{thm:gaugeConsistency}).
    \end{proof}
\end{theorem}

\begin{theorem}[Uniqueness of the normal-PEPS gauges up to a constant]%
    \label{thm:fundamentalTheorem_normalPEPS_gauge_unique}
    \lean{TNLean.PEPS.fundamentalTheorem_normalPEPS_gauge_unique}
    \leanok
    \uses{def:peps_normalPEPSBlockingHypotheses,
        def:peps_regionInjectivityDataPair, def:gaugeVertex}
    Let $A$ and $B$ be PEPS tensors on a finite simple graph with equal bond
    dimensions, the bond dimensions of the first tensor positive, satisfying
    the normal blocking hypotheses at the doubly injective assignment of the
    two tensors.  Suppose two families of invertible matrices, one on each
    edge, each realize the scalar-free gauge relation of
    Theorem~\ref{thm:fundamentalTheorem_normalPEPS_connected}: at every vertex
    the endpoint actions of the family transform $A_v$ into $B_v$.  Then the
    two families are proportional at every edge: $X'_e = c\cdot X_e$ for a
    nonzero scalar $c$ depending on the edge.  This is the uniqueness clause
    of the general normal-PEPS theorem stated after
    \cite[Theorem~3]{Molnar2018NormalPEPS}, read against the scalar-free
    relation; neither equality of the two states nor connectivity is needed.
    \begin{proof}
        \leanok
        \uses{thm:peps_gaugeConj_eq_of_coeffIdentities,
            lem:peps_injectiveRegion_union_pos}
        The inserted-edge coefficient is multilinear in the site tensors, one
        factor per site, so each scalar-free per-vertex relation makes every
        inserted coefficient of $B$ equal the corresponding inserted
        coefficient of the gauge-absorbed $A$: each family realizes the
        absorbed inserted-coefficient equality at every edge, with the roles
        of the two tensors exchanged.  On the first block of the edge's frame,
        whose blocked map and complement blocked map are injective for $A$ by
        the blocking hypotheses and union closure
        (Lemma~\ref{lem:peps_injectiveRegion_union_pos}), the region insertion
        transfer is determined by the absorbed equality
        (Theorem~\ref{thm:peps_gaugeConj_eq_of_coeffIdentities}), so the two
        orientation-adapted absorbing gauges induce the same conjugation and
        differ by a nonzero scalar; unwinding the orientation adaptation gives
        the proportionality of the gauges themselves.
    \end{proof}
\end{theorem}

\begin{theorem}[Fundamental Theorem for normal PEPS, vertex-injective case]%
    \label{thm:fundamentalTheorem_normalPEPS_of_vertexInjective}
    \lean{TNLean.PEPS.fundamentalTheorem_normalPEPS_of_vertexInjective}
    \leanok
    \uses{thm:fundamentalTheorem_PEPS, def:peps_normalPEPSBlockingHypotheses,
        def:GaugeEquivPEPS, def:IsVertexInjective}
    Suppose PEPS tensors $A$ and $B$ are vertex injective with positive bond
    dimensions on a connected graph, generate the same state, and $A$ satisfies
    the normal blocking hypotheses. Then their defining tensors are related by
    a local gauge.
    This is the vertex-injective specialization of the general normal-PEPS
    theorem: when each single tensor is already injective the gauge follows
    from the injective Fundamental Theorem, so the region blocking hypotheses
    are not used. The general region-injective theorem, in which only blocked
    regions are injective, remains open; the missing step is the
    region-level insertion-algebra isomorphism that produces the per-edge gauge
    after blocking.
\end{theorem}
\begin{proof}\leanok
    Vertex injectivity, positive bond dimensions, same state, and connectivity
    are exactly the hypotheses of the injective Fundamental Theorem, which
    supplies the local gauge.
\end{proof}

\begin{definition}[Oriented endpoint scalar]%
    \label{def:edgeScalarAt}
    \lean{TNLean.PEPS.edgeScalarAt}
    \leanok
    \uses{def:edgeGaugeAt}
    An edge-scalar family $c=(c_e)_{e\in E}$ determines the endpoint scalar
    $c_{v,e}$ by $c_{v,e}=c_e$ at the lower endpoint of $e$ and
    $c_{v,e}=c_e^{-1}$ at the upper endpoint.
\end{definition}

\begin{definition}[Vertex-balanced edge scalars]%
    \label{def:IsVertexBalanced}
    \lean{TNLean.PEPS.IsVertexBalanced}
    \leanok
    \uses{def:edgeScalarAt}
    An edge-scalar family is vertex-balanced when
    $\prod_{e\ni v} c_{v,e}=1$ for every vertex $v$.
\end{definition}
